% part: intuitionistic-logic
% chapter: soundness-completeness
% section: truth-lemma

\documentclass[../../../include/open-logic-section]{subfiles}

\begin{document}

\olfileid{int}{sc}{tru}

\olsection{सत्यता पूर्वप्रमेय}

पुढील पूर्वप्रमेय कॅनॉनिकल प्रतिमानातील
पूर्तीचा संबंध अभाज्य संच
$\Delta(\sigma)$ मध्ये कोणती
!!{formula}s !!{element}s आहेत
याच्याशी जोडते.

\begin{lem}\ollabel{lem:truth}
  $\Delta$ अभाज्य असेल, तर
  $\mSat{M(\Delta)}{!A}[\sigma]$
  तेव्हा आणि केवळ तेव्हाच
  $\Delta(\sigma) \Proves !A$.
\end{lem}

\begin{proof}
  $!A$ वर आगमन करू.
  \begin{enumerate}
    \item \indcase{!A}{\lfalse}{$\Delta(\sigma)$
      अभाज्य असल्याने सुसंगत आहे;
      म्हणून $\Delta(\sigma) \Proves/ \indfrm$.
      व्याख्येनुसार
      $\mSat/{M(\Delta)}{\indfrm}[\sigma]$.}
    \item \indcase{!A}{p}{$\mSat{{}}{}$ च्या
        व्याख्येनुसार
        $\mSat{M(\Delta)}{\indfrm}[\sigma]$
        तेव्हा आणि केवळ तेव्हाच
        $\sigma \in V(p)$;
        म्हणजे $\Delta(\sigma) \Proves \indfrm$.}
    \item \indcase!{!A}{\lnot !B}{}
    \item \indcase{!A}{!B \land
        !C}{$\mSat{M(\Delta)}{\indfrm}[\sigma]$
        तेव्हा आणि केवळ तेव्हाच
        $\mSat{M(\Delta)}{!B}[\sigma]$
        आणि $\mSat{M(\Delta)}{!C}[\sigma]$.
        आगमन-गृहीतकानुसार
        $\mSat{M(\Delta)}{!B}[\sigma]$
        तेव्हा आणि केवळ तेव्हाच
        $\Delta(\sigma) \Proves
        !B$; आणि~$!C$ साठीही तसेच.
        पण $\Delta(\sigma) \Proves !B$
        आणि $\Delta(\sigma) \Proves !C$
        हे तेव्हा आणि केवळ तेव्हाच,
        जेव्हा $\Delta(\sigma)
        \Proves \indfrm$.}
    \item \indcase{!A}{!B \lor
        !C}{$\mSat{M(\Delta)}{\indfrm}[\sigma]$
        तेव्हा आणि केवळ तेव्हाच
        $\mSat{M(\Delta)}{!B}[\sigma]$
        किंवा $\mSat{M(\Delta)}{!C}[\sigma]$.
        आगमन-गृहीतकानुसार हे
        तेव्हा आणि केवळ तेव्हाच,
        जेव्हा $\Delta(\sigma) \Proves !B$
        किंवा $\Delta(\sigma) \Proves !C$.
        आता हे पुन्हा
        $\Delta(\sigma) \Proves \indfrm$
        च्या तुल्य आहे, हे दाखवायचे आहे.
        डावीकडून उजवीकडील दिशा उघड आहे.
        उजवीकडून डावीकडील दिशा
        $\Delta(\sigma)$ अभाज्य
        असल्याने मिळते.}
    \item \indcase{!A}{!B \lif !C}{आधी
        डावीकडून उजवीकडील दिशेचे
        व्यत्यस्त रूप पाहू:
        $\Delta(\sigma) \Proves/ !B
        \lif !C$ असे समजा.
        मग $\Delta(\sigma) \cup \{!B\} \Proves/
        !C$ सुद्धा.
        कोणत्या तरी~$n$ साठी
        $\tuple{!B, !C}$ ही
        $\tuple{!B_n, !C_n}$ असल्याने
        $\Delta(\sigma.n) = (\Delta(\sigma) \cup
        \{!B\})^*$; आणि
        $\Delta(\sigma.n) \Proves !B$
        पण $\Delta(\sigma.n) \Proves/
        !C$. आगमन-गृहीतकानुसार
        $\mSat{M(\Delta)}{!B}[\sigma.n]$
        आणि $\mSat/{M(\Delta)}{!C}[\sigma.n]$.
        $R\sigma(\sigma.n)$ असल्यामुळे
        याचा अर्थ
        $\mSat/{M(\Delta)}{\indfrm}[\sigma]$.

        आता $\Delta(\sigma) \Proves !B \lif !C$
        असे समजा आणि $R\sigma\sigma'$ असू द्या.
        $\Delta(\sigma) \subseteq
        \Delta(\sigma')$ असल्याने
        $\Delta(\sigma') \Proves
        !B$ असेल, तर
        $\Delta(\sigma') \Proves !C$.
        दुसऱ्या शब्दांत, $R\sigma\sigma'$
        असलेल्या प्रत्येक $\sigma'$ साठी
        $\Delta(\sigma') \Proves/ !B$
        किंवा $\Delta(\sigma') \Proves
        !C$. आगमन-गृहीतकानुसार
        $R\sigma\sigma'$ असेल तेव्हा
        $\mSat/{M(\Delta)}{!B}[\sigma']$
        किंवा $\mSat{M(\Delta)}{!C}[\sigma']$;
        म्हणजे
        $\mSat{M(\Delta)}{\indfrm}[\sigma]$.}
  \end{enumerate}
\end{proof}

\end{document}
